%==============================================================================
% Appendix A: Notation and Conventions
%==============================================================================
\section{Notation and Conventions}
\label{app:A}

This appendix collects the notation conventions used in the body, the
shorthand for the four geometries, the terminology of the two-layer
structure, and the distinction between three 3-forms that are easily
conflated.  Since the formulas in the body are written in compact form, the
premises of the comparison are made explicit here.

\subsection{Geometries and Common Setup}
\label{app:A:geom}

\begin{itemize}
  \item $\Tthree$: flat three-torus.
  \item $\Nilthree$: Heisenberg-type nil geometry.
  \item $\Sthree$: round three-sphere.
  \item $\Solthree$: solv geometry.
\end{itemize}

In this paper ``the four geometries'' always refer to these four.  All
comparisons are carried out on the Euclidean product structure
\begin{equation}
  M^3 \times \Sone,
\end{equation}
with $L$ denoting the length of $\Sone$.  $r_0$ is the reference length
characterising the homogeneous background, and $R$ is the geometric scale
appearing in the structure constants of the three-dimensional spatial part.

\subsection{Two-Layer Structure and Response Vocabulary}
\label{app:A:layer}

\begin{itemize}
  \item \textbf{first layer}: response that can be defined on the bulk side
        without introducing a boundary.
  \item \textbf{second layer}: response defined in the
        boundary / interface / cobordism setting.
  \item \textbf{secondary reading}: an interpretive layer that re-reads, in
        another language, the already-fixed Euclidean response dictionary.
\end{itemize}

In the body we further use the following terms.

\begin{itemize}
  \item \textbf{allowed}: the channel is geometrically meaningful.
  \item \textbf{protected}: the channel is not automatically opened even
        under non-trivial perturbations.
  \item \textbf{activated}: a non-trivial response is raised by the
        structure of the geometry.
  \item \textbf{minimal}: there is no active family and the geometry sits
        on the benchmark side.
\end{itemize}

\subsection{Principal Symbols and Representative Values}
\label{app:A:symbols}

\begin{itemize}
  \item $\KMxw$: Maxwell baseline coefficient.
  \item $\SCZ$: Callias--Zanelli inflow action.
  \item $\etaAPS$: APS spectral invariant.
  \item $\Ntop$: frame-bundle-normalised torsional-charge entry.
  \item $C^2_\LC$: Levi--Civita Weyl scalar.
  \item spin-2 rigidity tag: scaffold-layer label
        (\texttt{trivial / non-rigid / rigid}).
  \item $\AKK$: KK anisotropy marker.
  \item AX torsion ansatz: $T_{ijk} = (2\eta/\rho)\,\epsilon_{ijk}$
        (DPPU AX-mode convention from paper~III (EC) \S3; the coefficient
        $2/\rho$ is fixed by this convention).
\end{itemize}

Representative benchmark values are listed below.  Derivation sketches and
verification scripts are referenced in Appendix~\ref{app:E}.

\begin{table}[H]
\centering
\resizebox{\linewidth}{!}{%
\begin{tabular}{lccl}
\toprule
Quantity & Value & Role & Derivation \\
\midrule
$\KMxw$ & $-L^2/(2 r_0^4)$ & universal bulk baseline & \S\ref{app:E:E2}, \S\ref{app:E:E5} \\
$\etaAPS(\Nilthree)$ & $+1/2$ & strict spectral core & \S\ref{app:E:E10} \\
$\Ntop(\Sthree)$ & $6 r_0^2$ & frame-bundle-norm.\ torsional-charge core & \S\ref{app:E:E9} \\
EC slice branch & present for $\Nilthree, \Solthree$ & spin-0 bulk branch & \S\ref{app:E:E5b} \\
$m^2(A_1) = m^2(A_2)$ in $\Solthree$ & $-L^2/(2 r_0^4)$ & biaxial KK splitting & \S\ref{app:E:E2}, \S\ref{app:E:E5} \\
$C^2_\LC(\Nilthree)$ & $4/(3 R^4)$ & scaffold datum & \S\ref{app:E:E7} \\
$C^2_\LC(\Solthree)$ & $16/(3 R^4)$ & scaffold datum & \S\ref{app:E:E7} \\
$k_{3D}^{\mathrm{tor\text{-}CS}}/k_{4D}^{\NY}$ & $1/2$ & NY--CS normalisation identity & \S\ref{app:E:E11} \\
$Q_{\mathrm{best}} = \eta(0)/\Delta k_{\mathrm{matter}}$ & $1$ & normalisation-free invariant & \S\ref{app:E:E12} \\
$k_q$ (formal) & $\in\mathbb{Z}$ & reduced CS level integrality & \S\ref{app:E:E13} \\
\bottomrule
\end{tabular}%
}
\caption{Representative benchmark values.}
\label{tab:appA_values}
\end{table}

\subsection{The Three 3-Forms and Their Connection}
\label{app:A:3forms}

We carefully separate the following 3-forms.

\begin{itemize}
  \item $\Ctor$: torsion-CS form.
  \item $\Csig$: spinor CS form.
  \item $\Cadj$: adjoint CS form.
\end{itemize}

The CZ inflow corresponds to $\Ctor$, and the APS spectral entry to $\Csig$.
In the body these must not be conflated.  Conceptually
\begin{equation}
  \mathrm{NY} = \dd\Ctor,
\end{equation}
and $\Cadj$ appears at intermediate stages of the APS argument.

\subsection{Geometric Scaffold Quantities}
\label{app:A:scaffold}

The following quantities, although not observable families themselves, are
used to record the distinctness of the background geometries.

\begin{itemize}
  \item $C^2_\LC$: magnitude of the Levi--Civita Weyl curvature.
  \item spin-2 rigidity tag: scaffold label summarising the zero pattern of
        the squash / shear Hessian.
  \item $\AKK / K^2$: magnitude of the KK anisotropy, with
        \begin{equation}
          \AKK := \sum_{i=0}^{2}(m_i^2 - \bar m^2)^2,
          \qquad
          \bar m^2 := \frac{1}{3}\sum_{i=0}^{2} m_i^2,
          \qquad
          K^2 := \KMxw^2.
        \end{equation}
\end{itemize}

This scaffold layer does not define an additional observable family, but it
is indispensable for distinguishing geometries such as $\Tthree$ and
$\Solthree$ that share the absence of a local CS core.  In particular,
$\Solthree$ is CS-inert yet carries an EC slice minimum and frame rigidity.

The benchmark values $\AKK / K^2$ can be computed directly from the mass
dictionaries:

\begin{itemize}
  \item $\Tthree$: $m_i^2 = 0$, so $\AKK / K^2 = 0$.
  \item $\Nilthree$ (uniaxial $m_2^2 = \KMxw$): $\bar m^2 = K/3$,
        $\AKK = 2(K/3)^2 + (2K/3)^2 = (2/3)K^2$.
  \item $\Sthree$ (triplet $m_i^2 = -2L^2/r_0^4$): all $i$ equal, so
        $\AKK / K^2 = 0$.
  \item $\Solthree$ (biaxial $m_1^2 = m_2^2 = \KMxw$): $\bar m^2 = 2K/3$,
        $\AKK = (2K/3)^2 + 2(K/3)^2 = (2/3)K^2$.
\end{itemize}

This symbolic confirmation is performed in
\texttt{proofs/kk\_higgsing.py}.
